\section{Experimental design}
\label{sec:experiments}

\subsection{Generated questions and partitioning}

The controlled dataset contains paired Japanese and English questions for four equally frequent families: addition, subtraction, multiplication, and a two-step operation.
The two-step family combines either addition followed by subtraction or multiplication followed by subtraction.
The generator computes integer reference answers directly from the sampled operands.
Names and objects come from small fixed vocabularies shared across partitions.
Each family has two main surface templates; separate templates provide the wording holdout.

The fixed generator seed is 20261108.
There are 1,024 training, 128 validation, 256 test and 128 wording-holdout questions, with equal numbers from each family in each partition.
Numeric operation-and-operand keys are disjoint across all four partitions.
Commutative operand pairs are sorted in these keys, including the first two operands of two-step questions.
This prevents a swapped multiplicative pair from appearing as the same arithmetic instance in both training and evaluation.
It does not make the arithmetic rules, individual numbers, names or objects unseen.
The wording holdout changes both the template and the numeric instance; it is not a paired paraphrase of the ordinary test set.

Before full runs, a small wiring pilot exposed four examples from each preliminary test partition.
That pilot revealed a generator key defect and a scoring distinction between numeric correctness and numeric-only formatting.
The corrected generator, both scoring rules and common decoding limits were fixed before full experiments.
This remains an exploratory evaluation rather than an untouched confirmatory test.
The protocol and generated JSONL files preserve the partition rules and the pilot amendments.

\subsection{External endpoint}

The external evaluation uses all 250 Japanese MGSM test questions from the official repository~\citep{shi2022mgsm}.
The loader pins repository revision \texttt{452a21ad3dae5668c06ceeac21ff073e1e40f9be}.
MGSM answers enter scoring only; neither its questions nor its answers are used for bridge fitting or checkpoint selection.
The bridge is therefore evaluated under a substantial change from its short synthetic training distribution.
Unlike the original MGSM study, this protocol supplies no few-shot chain-of-thought examples and asks for the final number only.
The results should not be compared directly with published chain-of-thought scores obtained under other prompting protocols.

\subsection{Fitting and comparisons}

The main bridge uses source layer 8 and prefix length 16, with independent training seeds 11, 23 and 37.
Layer 4, layer 16 and a four-token prefix are separate ablations at seed 11.
All fitted conditions use AdamW with learning rate $10^{-3}$, zero weight decay, gradient-norm clipping at 1.0, batch size 16 and ten epochs.
The saved checkpoint minimizes validation answer-token cross-entropy across the ten epochs.
No test score selects a checkpoint or the main layer.
The pretrained backbones run in BF16; the bridge parameters and cached pooled features use FP32.
The plotted training loss averages minibatch losses; validation loss is weighted by the number of supervised answer tokens.
Changing the bridge prefix length changes the target sequence length, while the shared affine map retains the same parameter count.

The learned textual soft prefix uses the same training examples, optimizer, epoch budget and checkpoint rule, at seed 11.
Its smaller parameter count and direct access to Japanese text remain material differences.
The untrained bridge uses the main architecture and seed without fitting.
The shuffled condition keeps the trained main bridge but permutes source feature sequences across evaluation questions, preserving an explicit mapping for each split.
The frozen textual systems receive either Japanese, the paired English template, or an actual PLaMo English translation.
The external endpoint has no paired-template English condition.

\subsection{Decoding, scoring and uncertainty}

Every answering condition uses greedy, free-running decoding with the same instruction to return the final number.
The cap is 128 new tokens for the generated tasks and 256 for MGSM.
PLaMo translation has caps of 128 and 512 tokens, respectively.
Per-item records include the generated text, emitted length, normal termination and cap-hit status.
The bridge path does not use generated translations or target answer tokens as inference inputs.

The originally specified strict exact match normalizes surrounding whitespace, special-token text, digit width, commas and a terminal period, and requires the complete remaining output to be the reference integer.
It therefore measures answer-format compliance together with correctness.
The secondary last-number exact match extracts the final numeric literal from the generated text and compares its exact decimal value with the integer reference.
For example, a correct arithmetic explanation ending in a number with a unit can pass the secondary score while failing the strict score.
This diagnostic is not a semantic verifier: an unfinished explanation can end in an operand rather than a final answer.
Cap-hit rates and raw predictions make this distinction inspectable.

Overall accuracy is the fraction of correct items; equal family counts make the synthetic overall score equal to the mean of the four family accuracies.
Selected differences use 10,000 paired bootstrap resamples of evaluation item IDs, with seed 20261108 and percentile 95\% intervals.
The same sampled IDs are used on both sides of each comparison.
These intervals are conditional on the fitted models and the particular evaluation set; they do not quantify all training variability.
The three main training seeds are reported separately and summarized without treating seeds as independent test items.
The ablations and their intervals are exploratory, with no multiple-comparison correction.
